\section{Обзор существующих решений}
\label{sec:Chapter2}

Ниже рассмотрены основные подходы к локализации БПЛА и оценивается степень их
соответствия требованиям, сформулированным в разделе~\ref{sec:Chapter1}:
работоспособность в отсутствие GNSS, сантиметровая точность на дистанциях в
десятки метров, низкая стоимость и независимость от дорогостоящей инфраструктуры.

\subsection{Спутниковая навигация (GNSS, RTK)}

Спутниковые системы (GPS, ГЛОНАСС и др.), в том числе дифференциальные и
RTK-решения, обеспечивают абсолютные координаты на открытой местности и при
хорошей видимости созвездия дают сантиметровую точность~\cite{kaplan2005,misra2006}.
Однако их точность и доступность существенно снижаются вблизи сооружений, внутри
помещений, в городской застройке и при загрязнённом радиоэфире~--- то есть именно
в тех условиях, для которых разрабатывается рассматриваемая система. Кроме того,
GNSS определяет координаты в глобальной системе отсчёта, тогда как для привязного
аппарата требуется положение \emph{относительно} платформы, которая сама может
перемещаться. Поэтому спутниковая навигация не удовлетворяет требованию работы в
условиях недоступности GNSS и плохо подходит для относительной локализации.

\subsection{Визуальные и визуально-инерциальные методы}

Оптический поток, визуальная одометрия и визуально-инерциальные системы (VIO)
способны обеспечивать высокую точность относительного позиционирования~\cite{scaramuzza2011}. Их
недостатки~--- зависимость от освещённости, текстуры сцены и динамики окружения, а
также значительные вычислительные затраты, что затруднительно для бортового
оборудования БПЛА с ограниченным энергобюджетом. В однородных или
слабоосвещённых сценах точность резко падает. Эти методы дополняют, но не заменяют
систему внешних маяков, когда требуется гарантированная точность относительно
платформы.

\subsection{Лидарная одометрия}

Лидарная одометрия отличается высокой информативностью и точностью построения
карты~\cite{zhang2014loam}, однако чувствительна к погодным условиям (туман, осадки, пыль), требует
дорогостоящего датчика и существенных вычислительных ресурсов, а также наличия
геометрических особенностей в сцене. Для лёгкого привязного аппарата масса,
стоимость и энергопотребление лидара являются серьёзными ограничениями.

\subsection{Системы захвата движения (motion capture)}

Инфраструктурные системы захвата движения обеспечивают субмиллиметровую точность~\cite{merriaux2017},
но применимы только в специально оборудованных зонах, требуют развёртывания сети
камер и калибровки и потому неприменимы в полевых условиях и противоречат
требованию независимости от дорогостоящей инфраструктуры.

\subsection{Радиочастотная дальнометрия (UWB)}

Сверхширокополосные (UWB) радиосистемы позволяют измерять расстояния по времени
распространения радиосигнала и широко используются для локализации внутри
помещений~\cite{gezici2005}. Их точность~--- порядка единиц--десятков сантиметров, что само по себе
близко к требуемой. Однако высокая скорость распространения радиоволн предъявляет
крайне жёсткие требования к разрешению измерения времени (пикосекундный масштаб):
ошибка в~$1$~нс соответствует~${\sim}30$~см дальности. При этом, как показано ниже,
радиоканал оказывается исключительно полезен не как измеритель дальности, а как
средство \emph{синхронизации} ультразвуковой подсистемы.

\subsection{Ультразвуковое позиционирование}

В ультразвуковых системах расстояние определяется по времени распространения
акустического сигнала. Низкая скорость звука (${\sim}340$~м/с против скорости
света) на пять с лишним порядков снижает требования к точности измерения времени:
для миллиметрового разрешения по дальности достаточно измерять время с точностью
порядка единиц микросекунд, что доступно недорогим микроконтроллерам.
Известны как схемы измерения времени пролёта (ToF), так и схемы разности времён
прихода (TDoA), для которых существуют эффективные методы решения задачи
локализации~\cite{smith1987,chan1994}. Ультразвуковые излучатели и приёмники
компактны, дёшевы и энергоэффективны. Ограничения метода~--- сравнительно малая
дальность и чувствительность к переотражениям~--- управляемы при работе в зоне
прямой видимости и временном кодировании сигнала.

\subsection{Вывод}

Сопоставление подходов с требованиями раздела~\ref{sec:Chapter1} показывает, что
наилучшим компромиссом для относительной локализации привязного БПЛА в
GNSS-недоступной среде является \textbf{ультразвуковое измерение дальности}
(низкая стоимость, малое энергопотребление, умеренные требования к измерению
времени, сантиметровая точность на нужных дистанциях). Ключевая трудность
ультразвуковой ToF-схемы~--- необходимость общей шкалы времени передатчика и
приёмника~--- решается \textbf{синхронизацией по UWB-радиоканалу}, который
обеспечивает точную временную привязку без жёстких требований к собственно
радиочастотной дальнометрии. Именно такая архитектура~--- ультразвуковая
дальнометрия с радиочастотной синхронизацией~--- положена в основу предлагаемого
решения (раздел~\ref{sec:Chapter3}).

\subsection{Существующие ультразвуковые системы локализации}
\label{sec:us-systems}

Сделанный выше выбор в пользу ультразвукового измерения дальности делает уместным обзор
существующих ультразвуковых систем локализации и сопоставление с ними предлагаемого решения.
Ультразвуковая локализация изучается и применяется не одно десятилетие; систематический
обзор технологий внутреннего позиционирования приведён, например, в~\cite{mautz2012}. Ниже
рассмотрены три показательные системы~--- две исследовательские и одна коммерческая.

\textbf{Active Bat} (AT\&T Laboratories Cambridge)~\cite{harter2002} использует подвижную
метку-излучатель и сетку ультразвуковых приёмников на потолке: по временам пролёта сигнала
до приёмников методом мультилатерации вычисляется трёхмерное положение с точностью около
$3$~см (в $95\%$ измерений ошибка не превышает $9$~см). Запуск измерения и временная
привязка выполняются по отдельному радиоканалу. Система точна, но требует плотной потолочной
инфраструктуры приёмников и ограничена объёмом помещения.

\textbf{Cricket} (MIT)~\cite{priyantha2000} строится на маяках, одновременно излучающих
радио- и ультразвуковой импульс; приёмник определяет расстояние до маяка по разности времён
прихода этих сигналов. Архитектура децентрализована~--- положение вычисляет сам приёмник;
точность составляет единицы сантиметров на дистанции в несколько метров, дальность одного
маяка~--- единицы метров. Система распространялась как относительно открытая платформа, но
давно не развивается.

\textbf{Коммерческие системы.} Среди серийных решений наиболее известна Marvelmind
Robotics~\cite{marvelmind}: стационарные ультразвуковые маяки, объединённые радиоканалом
диапазона $868/915$~МГц, с заявленной точностью порядка $\pm2$~см и рабочими дистанциями
между маяками до ${\sim}30$~м. Применяются и другие проприетарные системы позиционирования
реального времени (Sonitor, Hexamite). Коммерческие системы~--- законченные продукты с
многолетней эксплуатацией, однако их внутренняя реализация (алгоритмы дальнометрии,
обработки сигнала и синхронизации) закрыта.

\begin{table}[h!]
\centering
\small
\caption{Сравнение ультразвуковых систем локализации}
\label{tab:us-systems}
\begin{tabular}{lcccc}
\hline
Система & Точность & Дальность & Синхр. & Реализация \\
\hline
Active Bat    & ${\sim}3$~см       & помещение        & радио          & прототип          \\
Cricket       & ${\sim}2$~см       & единицы м        & РЧ\,+\,УЗ      & открытая (устар.) \\
Marvelmind    & $\pm2$~см          & до ${\sim}30$~м  & $868/915$~МГц  & закрытая          \\
Данная работа & ${\sim}5$~см (расч.) & до ${\sim}40$~м & UWB            & открытая          \\
\hline
\end{tabular}
\end{table}

Перечисленные системы ориентированы преимущественно на работу внутри помещений и опираются
на распределённую по среде инфраструктуру маяков или приёмников, достигая сантиметровой
точности (сводка~--- таблица~\ref{tab:us-systems}). Рассматриваемая в настоящей работе задача отличается сценарием: относительная
навигация привязного БПЛА над компактной (потенциально подвижной) платформой на дистанциях в
десятки метров, причём все якоря размещены на самой платформе, а не разнесены по среде.
Компактная база якорей повышает геометрическое ослабление точности (GDOP,
см.\ раздел~\ref{sec:Chapter3}), поэтому расчётная точность положения (${\sim}5$~см) несколько
уступает сантиметровому уровню инфраструктурных систем~--- это плата за отказ от стационарной
инфраструктуры. Вклад работы носит инженерный характер: открытая, воспроизводимая реализация
(схемотехника, прошивка, математическая модель) для указанного сценария, с синхронизацией
ультразвуковой подсистемы по UWB-радиоканалу. Превосходства над существующими системами по
точности работа не предполагает.
